% ============================================================
%  Chapter 13: The S-Entropy Calculus
%  Source: union/publication/support/unconstrained-subtask-recursion.tex
%  Parts I--II (S-scalar, receiver, axioms, triple equivalence structures)
% ============================================================

\epigraph{%
  Knowledge is distance from truth,\\
  measured on a bounded scale.%
}{}

%% ── §13.1 ─────────────────────────────────────────────────────────────────────
\section{Motivation: Why a Distinct Entropy?}
\label{sec:ch13-motivation}

Shannon entropy $H$ and thermodynamic entropy measure uncertainty, but neither
is indexed to a specific receiver's epistemic state relative to a specific
truth.  Shannon entropy is defined relative to a source distribution;
thermodynamic entropy is defined relative to a macrostate.  Neither gives the
\emph{distance} between a receiver's current knowledge and the correct answer to
a specific question.

The S-entropy calculus supplies this missing piece.  It is not a generalisation
of Shannon or Boltzmann entropy; it is a distinct axiomatic system that measures
distance to truth on a bounded closed interval, relative to a specified
receiver.  Every quantity that appears in the calculus---the floor, the
catalytic power, the triple---is receiver-relative.  The global truth is
asserted by the global S-value; what happens at the level of subtasks is
unconstrained.

The calculus is presented as pure mathematics.  Its application to mass
spectrometry follows in the subsequent chapters as a special case: the
S-entropy triple $(\Sk, \St, \Se)$ instantiates the abstract triple
$(S_k, S_t, S_e)$ of the source paper for the specific receiver defined by a
mass spectrometric apparatus.

%% ── §13.2 ─────────────────────────────────────────────────────────────────────
\section{The S-Scale}
\label{sec:ch13-scale}

\begin{definition}[S-scale]\label{def:sscale-mono}
The \emph{S-scale} is the closed real interval
\[
  \mathbb{S} := [0,100] \subset \mathbb{R}.
\]
We write $\mathbb{S}^{*} := (0,100]$ for the open-below S-scale and
$\partial\mathbb{S} := \{0,100\}$ for the two boundary points: $S=0$ denotes
exact alignment with truth (perfect knowledge), and $S=100$ denotes maximal
misalignment (complete ignorance).
\end{definition}

The bounds are not arbitrary.  $S = 100$ corresponds to the maximum-entropy
state of a receiver whose knowledge is entirely uncorrelated with the truth.
$S = 0$ is the limit of perfect knowledge---the Floor Theorem
(Chapter~\ref{ch:floor}) will show that no bounded receiver can attain it.
The interval $[0,100]$ is chosen so that S-values are dimensionless and
universally comparable across receivers.

\begin{remark}[S-scale versus information-theoretic quantities]
The S-scale is \emph{not} a logarithmic measure.  It is a direct distance
measure on the space of candidates relative to a truth, normalised to the
compact interval $[0,100]$.  The logarithmic structure of Shannon entropy is
replaced here by the floor positivity axiom (Axiom~(S3) below), which encodes
the unavoidable residual distance characteristic of bounded receivers.
\end{remark}

%% ── §13.3 ─────────────────────────────────────────────────────────────────────
\section{Candidate Space, Truth Space, and the Receiver Quadruple}
\label{sec:ch13-receiver}

\begin{definition}[Candidate space and truth space]\label{def:cand-truth-mono}
A \emph{candidate space} $\mathcal{X}$ is a non-empty set whose elements
represent admissible candidate solutions to a problem.  A
\emph{truth space} $\mathcal{X}^{*} \subseteq \mathcal{X}$ is a designated
non-empty subset of correct solutions.  We require $|\mathcal{X}^{*}| \geq 1$;
the case $|\mathcal{X}^{*}| > 1$ allows for solutions that are operationally
indistinguishable.
\end{definition}

\begin{definition}[Receiver]\label{def:receiver-mono}
A \emph{receiver} $\mathcal{R}$ over $(\mathcal{X}, \mathcal{X}^{*})$ is a
quadruple
\[
  \mathcal{R} = (\Sigma_{\mathcal{R}},\; \Phi_{\mathcal{R}},\;
                 K_{\mathcal{R}},\; \mathrm{dec}_{\mathcal{R}})
\]
where:
\begin{enumerate}[(i)]
  \item $\Sigma_{\mathcal{R}}$ is a non-empty set, the \emph{signal space} of
    $\mathcal{R}$;
  \item $\Phi_{\mathcal{R}} : \Sigma_{\mathcal{R}} \to K_{\mathcal{R}}$
    is the \emph{decoder} function;
  \item $K_{\mathcal{R}}$ is the \emph{knowledge framework} of $\mathcal{R}$,
    equipped with a fixed set of operations
    $\{\circ_j\}_{j \in J_{\mathcal{R}}}$;
  \item $\mathrm{dec}_{\mathcal{R}} : K_{\mathcal{R}} \to \mathcal{X}$ is the
    \emph{candidate-projection}, associating each knowledge state with a
    candidate.
\end{enumerate}
The receiver is \emph{bounded} if $|K_{\mathcal{R}}|$ is a cardinal strictly
less than that of $\mathcal{X}$ when $|\mathcal{X}|$ is infinite.
\end{definition}

The four components of the receiver have precise physical interpretations in
the mass spectrometric setting.  $\Sigma_{\mathcal{R}}$ is the space of raw
detector signals; $\Phi_{\mathcal{R}}$ is the signal-processing pipeline that
maps raw signals to knowledge states (peak lists, identified features);
$K_{\mathcal{R}}$ is the spectral library or internal reference database; and
$\mathrm{dec}_{\mathcal{R}}$ maps each library entry to a molecular candidate.
The boundedness of the receiver is the statement that the library has strictly
fewer entries than the space of all possible molecular candidates.

\begin{remark}[The Axiom of Boundedness and receivers]
Every persistent physical instrument is a bounded receiver.  The signal space
$\Sigma_{\mathcal{R}}$ satisfies the Bounded Phase Space Law
(Axiom~\ref{ax:bpsl}): it has finite Liouville measure, is measure-preserving,
and admits a hierarchical partition.  The boundedness clause of
Definition~\ref{def:receiver-mono} is therefore a consequence of
Axiom~\ref{ax:bpsl} applied to the instrument as a dynamical system.
\end{remark}

%% ── §13.4 ─────────────────────────────────────────────────────────────────────
\section{The S-Functional: Four Axioms}
\label{sec:ch13-functional}

\begin{definition}[S-functional]\label{def:sfunc-mono}
An \emph{S-functional} on $(\mathcal{X}, \mathcal{X}^{*})$ is a map
\[
  S : \mathcal{R}\text{-class} \times \mathcal{X} \times \mathcal{X}^{*}
    \;\longrightarrow\; \mathbb{S},
  \qquad
  (\mathcal{R}, x, x^{*}) \;\mapsto\; S_{\mathcal{R}}(x, x^{*}),
\]
satisfying the four axioms (S1)--(S4) below.
\end{definition}

\begin{theorem}[(S1) Non-negativity]\label{ax:S1-mono}
For every receiver $\mathcal{R}$, candidate $x \in \mathcal{X}$, and truth
$x^{*} \in \mathcal{X}^{*}$:
\[
  S_{\mathcal{R}}(x, x^{*}) \;\geq\; 0.
\]
\end{theorem}

\begin{theorem}[(S2) Boundedness]\label{ax:S2-mono}
For every receiver $\mathcal{R}$, candidate $x \in \mathcal{X}$, and truth
$x^{*} \in \mathcal{X}^{*}$:
\[
  S_{\mathcal{R}}(x, x^{*}) \;\leq\; 100.
\]
\end{theorem}

\begin{theorem}[(S3) Self-truth lower bound (Floor Axiom)]\label{ax:S3-mono}
There exists a function $S_{\flat} : \mathcal{R}\text{-class} \to \mathbb{S}^{*}$,
called the \emph{floor function}, with $S_{\flat}(\mathcal{R}) > 0$ for every
receiver $\mathcal{R}$, such that
\[
  \inf_{x \in \mathcal{X}}\; S_{\mathcal{R}}(x, x^{*})
  \;=\; S_{\flat}(\mathcal{R})
\]
for every $x^{*} \in \mathcal{X}^{*}$.
\end{theorem}

\begin{theorem}[(S4) Receiver-decoded coherence]\label{ax:S4-mono}
For every receiver $\mathcal{R}$, the partial map
$x \mapsto S_{\mathcal{R}}(x, x^{*})$ is
$\mathrm{dec}_{\mathcal{R}}$-coherent: there exists a non-decreasing function
\[
  \delta_{\mathcal{R}} :
  K_{\mathcal{R}} \times K_{\mathcal{R}} \to \mathbb{S}
\]
with $\delta_{\mathcal{R}}(\sigma, \sigma) = S_{\flat}(\mathcal{R})$ for all
$\sigma$, such that
\[
  S_{\mathcal{R}}(x, x^{*})
  \;=\; \delta_{\mathcal{R}}\!\bigl(
    \Phi_{\mathcal{R}}(\sigma_x),\;
    \Phi_{\mathcal{R}}(\sigma_{x^{*}})\bigr)
\]
for some $\sigma_x \in \Sigma_{\mathcal{R}}$ with
$\mathrm{dec}_{\mathcal{R}}(\Phi_{\mathcal{R}}(\sigma_x)) = x$,
and similarly for $\sigma_{x^{*}}$.
\end{theorem}

These four axioms are the complete specification of the S-functional.  Axioms
(S1) and (S2) fix the codomain to the S-scale.  Axiom (S3) is the structural
heart of the theory: it asserts that the infimum of all S-values attainable
by the receiver is a strictly positive constant $S_{\flat}(\mathcal{R})$
depending only on the receiver's internal structure, not on the candidate or
the truth.  Axiom (S4) requires that the S-value be computable from the
receiver's internal decoding of signals.

\begin{remark}[Axiom (S3) and bounded cognition]
Axiom~(S3) expresses the boundedness of cognition: no receiver can attain
perfect alignment with the truth.  The constant $S_{\flat}(\mathcal{R})$ is
the \emph{smallest knowable error} of $\mathcal{R}$.  The Floor Theorem of
Chapter~\ref{ch:floor} derives this positivity rigorously from the cardinality
structure of bounded receivers; Axiom~(S3) here merely stipulates the property
as an axiom, leaving the derivation to the theorem.
\end{remark}

\begin{lemma}[Image of $S_{\mathcal{R}}$]\label{lem:image-S-mono}
For each receiver $\mathcal{R}$, the image of $S_{\mathcal{R}}$ is contained
in $[S_{\flat}(\mathcal{R}), 100] \subseteq \mathbb{S}^{*}$.
\end{lemma}

\begin{proof}
Immediate from Axioms~(S1)--(S3): by (S1) and (S2) the image lies in $[0,100]$,
and by (S3) no value below $S_{\flat}(\mathcal{R})$ is attained.
\end{proof}

%% ── §13.5 ─────────────────────────────────────────────────────────────────────
\section{Symmetric S-Functionals}
\label{sec:ch13-symmetric}

\begin{definition}[Symmetric S-functional]\label{def:sym-S-mono}
An S-functional $S$ is \emph{symmetric} if for every receiver $\mathcal{R}$
and every $x, y \in \mathcal{X} \cup \mathcal{X}^{*}$:
\begin{align}
  S_{\mathcal{R}}(x, y) &= S_{\mathcal{R}}(y, x)
  \quad\text{(symmetry),} \label{eq:sym}\\
  S_{\mathcal{R}}(x, z) &\leq S_{\mathcal{R}}(x, y)
    + S_{\mathcal{R}}(y, z)
  \quad\text{(triangle inequality, modulo floor).} \label{eq:tri}
\end{align}
\end{definition}

\begin{proposition}[Floor on symmetric S-functionals]\label{prop:sym-floor-mono}
If $S$ is symmetric and the triangle inequality~\eqref{eq:tri} holds modulo
floor $S_{\flat}(\mathcal{R})$, then for every distinct
$x, y \in \mathcal{X}$:
\[
  S_{\mathcal{R}}(x, y) \;\geq\; S_{\flat}(\mathcal{R}).
\]
\end{proposition}

\begin{proof}
Apply the triangle inequality with $z = x$:
$S_{\mathcal{R}}(x, x) \leq S_{\mathcal{R}}(x, y) + S_{\mathcal{R}}(y, x)
= 2\,S_{\mathcal{R}}(x, y)$ by symmetry.
Since $S_{\mathcal{R}}(x, x) \geq S_{\flat}(\mathcal{R})$ by the diagonal
floor (Axiom~(S4)), we obtain
$S_{\mathcal{R}}(x, y) \geq S_{\flat}(\mathcal{R})/2$.
By induction on the hop count between $x$ and $y$ in any finite cover, the
bound tightens to $S_{\flat}(\mathcal{R})$.
\end{proof}

\begin{remark}[Semimetric structure]
Symmetric S-functionals satisfying both conditions of
Definition~\ref{def:sym-S-mono} form a \emph{semimetric} on $\mathcal{X}$
with offset $S_{\flat}$.  The calculus is developed for the general (possibly
asymmetric) case; symmetric variants are the common case arising from physical
instruments.
\end{remark}

%% ── §13.6 ─────────────────────────────────────────────────────────────────────
\section{Bounded Cognition and the Cognitive Residue}
\label{sec:ch13-cognition}

\begin{definition}[Cognitive capacity]\label{def:cog-cap-mono}
The \emph{cognitive capacity} of a receiver $\mathcal{R}$ is the cardinal
$|K_{\mathcal{R}}|$.  We say $\mathcal{R}$ is \emph{finite} if
$|K_{\mathcal{R}}|$ is finite, and \emph{denumerably bounded} if
$|K_{\mathcal{R}}| \leq \aleph_0$.
\end{definition}

\begin{definition}[Total information space]\label{def:total-info-mono}
The \emph{total information space} of a problem $(\mathcal{X}, \mathcal{X}^{*})$
is a set $\mathcal{U}_{\mathrm{tot}}$ with
$\mathcal{X} \subseteq \mathcal{U}_{\mathrm{tot}}$, equipped with a cardinality
$|\mathcal{U}_{\mathrm{tot}}|$ that may exceed $|\mathcal{X}|$.
\end{definition}

\begin{definition}[Cognitive residue]\label{def:residue-mono}
The \emph{cognitive residue} of a receiver $\mathcal{R}$ relative to total
information space $\mathcal{U}_{\mathrm{tot}}$ is
\[
  \mathcal{R}\text{-residue}(\mathcal{R})
  \;:=\; \mathcal{U}_{\mathrm{tot}} \setminus
    \mathrm{dec}_{\mathcal{R}}\!\bigl(
      \Phi_{\mathcal{R}}(\Sigma_{\mathcal{R}})\bigr).
\]
\end{definition}

\begin{theorem}[Residue persistence]\label{thm:residue-mono}
Let $\mathcal{R}$ be a receiver with cognitive capacity
$|K_{\mathcal{R}}| < |\mathcal{U}_{\mathrm{tot}}|$.  Then the cognitive
residue satisfies
\[
  |\mathcal{R}\text{-residue}(\mathcal{R})|
  \;=\; |\mathcal{U}_{\mathrm{tot}}|.
\]
\end{theorem}

\begin{proof}
Since $\mathrm{dec}_{\mathcal{R}} : K_{\mathcal{R}} \to \mathcal{X}$ is a
function,
$|\mathrm{dec}_{\mathcal{R}}(\Phi_{\mathcal{R}}(\Sigma_{\mathcal{R}}))| \leq
|K_{\mathcal{R}}|$.
Therefore
$|\mathcal{R}\text{-residue}(\mathcal{R})| \geq
|\mathcal{U}_{\mathrm{tot}}| - |K_{\mathcal{R}}|$.
By the cardinal arithmetic of subtraction in infinite cardinals, when
$|K_{\mathcal{R}}| < |\mathcal{U}_{\mathrm{tot}}|$ this difference equals
$|\mathcal{U}_{\mathrm{tot}}|$.
\end{proof}

The residue persistence theorem is the cardinality reason that no bounded
receiver can attain $S = 0$: there always remain unrepresented elements of
the total information space, and these elements contribute unavoidable distance
to truth.  Chapter~\ref{ch:floor} makes this argument into the Floor Theorem.

%% ── §13.7 ─────────────────────────────────────────────────────────────────────
\section{The S-Entropy Triple for Mass Spectrometry}
\label{sec:ch13-triple}

The abstract calculus produces a triple of S-values whenever the receiver's
knowledge framework separates into three relatively independent components.
For a mass spectrometric receiver, the three natural components are:

\begin{enumerate}[(i)]
  \item \textbf{Oscillatory S-entropy} $\Sk$: the S-value for the question
    ``what is the molecular structure?''\ answered through spectral matching
    and fragmentation pattern recognition.  $\Sk$ measures how far the
    receiver's current spectral assignment is from the true molecular identity.

  \item \textbf{Categorical S-entropy} $\St$: the S-value for the question
    ``to which category does this ion belong?''\ resolved via database lookup,
    spectral library search, and retention-time prediction.  $\St$ measures
    distance from the correct categorical assignment in the receiver's
    classification hierarchy.

  \item \textbf{Partition S-entropy} $\Se$: the S-value for the question
    ``which partition cell does this ion occupy?''\ at the instrument's current
    resolution depth.  $\Se$ measures how far the receiver's finest achievable
    $m/z$ partitioning is from the true mass coordinate of the ion.
\end{enumerate}

\begin{definition}[Mass spectrometric S-triple]\label{def:ms-triple}
The \emph{mass spectrometric S-entropy triple} is the ordered triple
\[
  (\Sk,\; \St,\; \Se)
  \;\in\; \mathbb{S}^3
\]
where each component is the S-value of the respective question for the given
mass spectrometric receiver $\mathcal{R}_{\mathrm{MS}}$ at the current
knowledge state $K_{\mathcal{R}}$.
\end{definition}

Each component is independently bounded in $[S_{\flat}(\mathcal{R}), 100]$ by
Lemma~\ref{lem:image-S-mono}.  Their joint variation maps the instrument's
epistemic state: as the instrument processes more scans and refines its
assignments, all three components decrease toward the floor, but none reaches
zero.  The Triple Equivalence theorem of Chapter~\ref{ch:triple} will show
that $\Sk = \St = \Se$ at any fixed depth of the partition hierarchy---the
three entropies are the same quantity measured in three distinct observational
languages.

\begin{remark}[The triple and the Unconstrained Subtask Theorem]
A key structural fact, proved in Chapter~\ref{ch:floor} as the Unconstrained
Subtask Theorem, is that the global S-value $S_{\mathcal{R}}(\xi, x^{*})$ of
an expression $\xi$ does not constrain the local S-values of its subtasks.
In particular: the global spectral assignment task has a fixed S-value
$S^{*}$, but the individual sub-tasks (peak picking, isotope correction,
database lookup, retention-time matching) may each have \emph{any} local
S-value in $[S_{\flat}, 100]$---including $100$, meaning locally maximal
ignorance---and the global answer can still be correct.  This is the formal
justification for the GPU multi-pass architecture: each shader pass operates
on its own local S-value independently, and the composition yields the correct
global S-value.
\end{remark}

%% ── §13.8 ─────────────────────────────────────────────────────────────────────
\section{Three Algebraic Structures Underlying the Triple}
\label{sec:ch13-structures}

Before stating the Triple Equivalence (Chapter~\ref{ch:triple}), we record the
three algebraic structures that the categories $\Osc$, $\Cat$, and $\Part$
formalise.  These definitions are the abstract counterparts of the three
S-entropy components.

\begin{definition}[Oscillatory structure]\label{def:osc-mono}
An \emph{oscillatory structure} is a quadruple
$\mathfrak{O} = (M, \omega, \Phi, \mu)$
where:
\begin{enumerate}[(i)]
  \item $M$ is a non-empty set, the \emph{mode set};
  \item $\omega : M \to \mathbb{R}_{\geq 0}$ is the \emph{frequency
    assignment};
  \item $\Phi : M \times \mathbb{R} \to S^1$ is the \emph{phase map},
    satisfying $\Phi(m, t+s) = \Phi(m,t) \cdot e^{2\pi i\,\omega(m)s}$;
  \item $\mu$ is a probability measure on $M$ encoding mode weights.
\end{enumerate}
Morphisms of oscillatory structures preserve frequencies and phase
compatibility; the resulting category is written $\mathbf{Osc}$.
\end{definition}

\begin{definition}[Categorical structure]\label{def:cat-mono}
A \emph{categorical structure} is a quadruple
$\mathfrak{E} = (C, \prec, \lambda, \nu)$
where:
\begin{enumerate}[(i)]
  \item $C$ is a non-empty set, the \emph{cell set};
  \item $\prec$ is a strict partial order on $C$ with the property that
    $\{c' : c' \prec c\}$ is finite for each $c \in C$
    (the \emph{refinement order});
  \item $\lambda : C \to \Lambda$ is the \emph{labelling}, where $\Lambda$ is
    a set of label tuples;
  \item $\nu$ is a probability measure on $C$ encoding cell weights.
\end{enumerate}
Morphisms preserve $\prec$ and $\lambda$; the resulting category is
written $\mathbf{Cat}$.
\end{definition}

\begin{definition}[Partition structure]\label{def:part-mono}
A \emph{partition structure} is a triple
$\mathfrak{P} = (\Omega, \mathscr{P}, \pi)$
where:
\begin{enumerate}[(i)]
  \item $\Omega$ is a non-empty measurable space with $\sigma$-algebra
    $\mathscr{B}$;
  \item $\mathscr{P} = \{P_n\}_{n \in \mathbb{N}}$ is a sequence of finite
    measurable partitions of $\Omega$ with $P_{n+1}$ refining $P_n$;
  \item $\pi$ is a probability measure on $\Omega$ compatible with each $P_n$
    in the sense that $\pi(\partial p) = 0$ for each $p \in P_n$.
\end{enumerate}
Morphisms preserve refinement and the measure; the resulting category is
written $\mathbf{Part}$.
\end{definition}

The three structures encode the same information about the underlying
probability space in three different observational languages.  This is the
content of the Triple Equivalence: the choice among
$\mathbf{Osc}$, $\mathbf{Cat}$, $\mathbf{Part}$ is \emph{computational}, not
informational.  A computation may be performed in whichever representation is
cheapest.

\begin{remark}[Connection to Axiom~\ref{ax:bpsl}]
The partition structure $\mathfrak{P} = (\Omega, \mathscr{P}, \pi)$ is
precisely the structure guaranteed by clause~(iii) of the Bounded Phase Space
Law (Axiom~\ref{ax:bpsl}): every persistent physical system has a hierarchical
partition $\mathcal{P}_1 \prec \mathcal{P}_2 \prec \cdots$ of its phase space.
The S-entropy calculus is therefore grounded directly in the Axiom: the
partition structure required for bounded phase space is the same partition
structure required for S-entropy.
\end{remark}

%% ── Chapter Summary ───────────────────────────────────────────────────────────
\bigskip
\begin{tcolorbox}[colback=crownblue!5, colframe=crownblue!60!black,
                  title={\textbf{Chapter 13 Summary}}]
\begin{itemize}[noitemsep]
  \item The S-scale is the bounded interval $\mathbb{S} = [0,100]$; $S=0$
    denotes perfect knowledge, $S=100$ complete ignorance, and the Floor
    Theorem (Chapter~\ref{ch:floor}) establishes that $S=0$ is
    unattainable by any bounded receiver.
  \item A receiver is a quadruple
    $\mathcal{R} = (\Sigma_{\mathcal{R}}, \Phi_{\mathcal{R}},
    K_{\mathcal{R}}, \mathrm{dec}_{\mathcal{R}})$; a mass spectrometer is a
    bounded receiver in this sense.
  \item The S-functional satisfies four axioms: (S1)~non-negativity,
    (S2)~boundedness, (S3)~self-truth lower bound (floor positivity),
    and (S4)~receiver-decoded coherence.
  \item The cognitive residue of a bounded receiver has the same cardinality as
    the total information space; this is the cardinality source of the
    unavoidable floor $S_{\flat}(\mathcal{R}) > 0$.
  \item The mass spectrometric S-entropy triple is $(\Sk, \St, \Se)$,
    corresponding to oscillatory (structural), categorical (classification),
    and partition (resolution) distances to truth.
  \item The three algebraic structures $\mathbf{Osc}$, $\mathbf{Cat}$,
    $\mathbf{Part}$ are defined; their mutual equivalence is the Triple
    Equivalence of Chapter~\ref{ch:triple}.
\end{itemize}
\end{tcolorbox}
